\section{КОНТЕЙНЕРИЗАЦИЯ, ТЕСТИРОВАНИЕ И РАЗВЕРТЫВАНИЕ}

\subsection{Спецификация контейнеров (Docker)}
Для каждого микросервиса создан легковесный образ на базе Python. Это гарантирует, что система будет работать одинаково на любой машине, где установлен Docker.

\subsection{Сценарии верификации работоспособности}
Проверка системы осуществляется через последовательность действий:
\begin{enumerate}
    \item Сборка и запуск: \texttt{docker-compose up --build}.
    \item Создание объекта: запрос \texttt{POST /api/invoices}.
    \item Чтение объекта: запрос \texttt{GET /api/invoices/1}.
    \item Проверка списка: запрос \texttt{GET /api/invoices}.
\end{enumerate}

\subsection{Автоматизация сборки и тестирования}
Сборка проекта и запуск тестов автоматизированы с помощью \texttt{Makefile}. Тесты проверяют наличие всех необходимых файлов и корректность указания \texttt{project\_code} (invoices-s13) в документации.

\subsection{Ограничения и потенциал развития}
На текущем этапе система является функциональным прототипом. В качестве дальнейшего развития планируется внедрение персистентной БД (PostgreSQL) и системы аутентификации на базе JWT-токенов.
